\uuid{m14q}
\titre{Bijection}
\auteur{bodin ; hanani ; mzari}
\organisation{exo7}
\licence{CC BY-NC-SA 4.0}
\datecreate{2020-01-20}
\niveau{L1}
\module{Analyse}
\chapitre{Continuité, limite et étude de fonctions réelles}
\sousChapitre{Etude de fonctions}
\difficulte{2}
\theme{Maximum, bijection, bijection}
\exo7qcmid{lille-1.continuite.028}

\contenu{
\begin{question}
\qclassification{123.04}
\qtags{motcle=bijection}
Soit $f : [a,b] \to [c,d]$ continue avec $a < b$, $c < d$, $f(a)=c$, $f(b)=d$. Quelles propriétés impliquent $f$ bijective ?
\begin{answers}
    \good{$f$ injective.}

    \bad{$f$ surjective.}

    \bad{$f$ croissante.}

    \good{$f$ strictement croissante.}
\end{answers}
\begin{explanations}
Comme $f(a)=c$, $f(b)=d$, alors par le théorème des valeurs intermédiaires, toute valeur entre $c$ et $d$ est atteinte, autrement dit $f$ est surjective.
Si en plus $f$ est injective (ce qui est le cas si $f$ strictement croissante) alors $f$ sera bijective.
\end{explanations}
\end{question}
}
